% Pipeline schematic (TikZ). Wrapped in \resizebox so it always fits the
% text width in both IEEEtran and the CVF/wacv two-column layout.
\begin{figure*}[t]
\centering
\resizebox{\textwidth}{!}{%
\begin{tikzpicture}[
  font=\footnotesize,
  >=Stealth,
  box/.style={draw, rounded corners, minimum height=9mm, align=center, inner sep=3pt},
  data/.style={box, fill=gray!10},
  frozen/.style={box, fill=blue!8},
  train/.style={box, fill=violet!12},
  loss/.style={box, fill=red!10}
]
% camera branch (top)
\node[data]   (img)   at (0,1.3)   {Camera\\image};
\node[frozen] (dino)  at (2.3,1.3)  {DINOv2\\ViT-B/14\\(frozen)};
\node[data]   (feat)  at (4.7,1.3)  {patch feat.\\$16{\times}16{\times}768$};
% lidar branch (bottom)
\node[data]   (lidar) at (0,-1.3)  {LiDAR\\points};
\node[data]   (vox)   at (2.3,-1.3) {voxelize\\$128^2{\times}32$};
% fused spine
\node[train]  (fuse)  at (7.2,0)   {voxel input\\$\mathbf{O}\oplus\mathbf{F}^v$};
\node[train]  (mask)  at (9.3,0)   {range-aware\\mask $\rho{=}0.85$};
\node[train]  (enc)   at (11.3,0)  {3D CNN\\encoder};
\node[train]  (bott)  at (13.0,0)  {bottleneck\\$256{\times}32^2{\times}8$};
\node[train]  (dec)   at (14.9,0)  {3D CNN\\decoder};
% heads + losses
\node[box,fill=violet!12] (occh) at (16.9,0.95) {head:\\occ.\,$\hat{\mathbf{O}}$};
\node[box,fill=violet!12] (fth)  at (16.9,-0.95) {head:\\feat.\,$\hat{\mathbf{F}}^v$};
\node[loss] (locc) at (18.9,0.95) {$\mathcal{L}_{\text{occ}}$\\BCE};
\node[loss] (lfeat) at (18.9,-0.95) {$\mathcal{L}_{\text{feat}}$\\MSE};
% arrows
\draw[->] (img)  -- (dino);
\draw[->] (dino) -- (feat);
\draw[->] (lidar)-- (vox);
\draw[->] (feat)  -- (fuse);
\draw[->] (vox)   -- (fuse);
\draw[->] (fuse) -- (mask);
\draw[->] (mask) -- (enc);
\draw[->] (enc)  -- (bott);
\draw[->] (bott) -- (dec);
\draw[->] (dec) -- (occh);
\draw[->] (dec) -- (fth);
\draw[->] (occh) -- (locc);
\draw[->] (fth)  -- (lfeat);
\end{tikzpicture}%
}
\caption{\textbf{OMU-MAE pretraining pipeline.} A paired RGB image (top)
is processed by a \emph{frozen} DINOv2 ViT-B/14 encoder; per-patch
features are projected via the camera--LiDAR calibration and
mean-aggregated per voxel. The LiDAR scan (bottom) is voxelized at
$0.4$\,m into a $128{\times}128{\times}32$ grid. The concatenated
cross-modal voxel input $\mathbf{O}\oplus\mathbf{F}^v$ is range-aware
masked ($\rho{=}0.85$), encoded by a 3D CNN, and decoded back to two
heads that reconstruct occupancy ($\mathcal{L}_{\text{occ}}$, BCE) and
DINOv2 features ($\mathcal{L}_{\text{feat}}$, MSE) at masked positions.}
\label{fig:pipeline}
\end{figure*}
